\documentclass[10pt,a4paper]{amsart}
\usepackage{amsmath,amssymb,amsthm,graphicx}
\usepackage{verbatim}
\newenvironment{reference}{\comment}{\endcomment}
%\newenvironment{reference}{}{}
\newenvironment{slogan}{\comment}{\endcomment}
\newenvironment{history}{\comment}{\endcomment}

% For commutative diagrams we use Xy-pic
\usepackage[all]{xy}

% We use 2cell for 2-commutative diagrams.
\xyoption{2cell}
% Theorem environments.
%
\theoremstyle{plain}
\newtheorem{theorem}[subsection]{Theorem}
\newtheorem{proposition}[subsection]{Proposition}
\newtheorem{lemma}[subsection]{Lemma}

\theoremstyle{definition}
\newtheorem{definition}[subsection]{Definition}
\newtheorem{example}[subsection]{Example}
\newtheorem{exercise}[subsection]{Exercise}
\newtheorem{situation}[subsection]{Situation}

\theoremstyle{remark}
\newtheorem{remark}[subsection]{Remark}
\newtheorem{remarks}[subsection]{Remarks}

\numberwithin{equation}{subsection}

% Macros
%
\def\lim{\mathop{\mathrm{lim}}\nolimits}
\def\colim{\mathop{\mathrm{colim}}\nolimits}
\def\Spec{\mathop{\mathrm{Spec}}}
\def\Hom{\mathop{\mathrm{Hom}}\nolimits}
\def\Ext{\mathop{\mathrm{Ext}}\nolimits}
\def\SheafHom{\mathop{\mathcal{H}\!\mathit{om}}\nolimits}
\def\SheafExt{\mathop{\mathcal{E}\!\mathit{xt}}\nolimits}
\def\Sch{\mathit{Sch}}
\def\Mor{\mathop{\mathrm{Mor}}\nolimits}
\def\Ob{\mathop{\mathrm{Ob}}\nolimits}
\def\Sh{\mathop{\mathit{Sh}}\nolimits}
\def\NL{\mathop{N\!L}\nolimits}
\def\CH{\mathop{\mathrm{CH}}\nolimits}
\def\proetale{{pro\text{-}\acute{e}tale}}
\def\etale{{\acute{e}tale}}
\def\QCoh{\mathit{QCoh}}
\def\Ker{\mathop{\mathrm{Ker}}}
\def\Im{\mathop{\mathrm{Im}}}
\def\Coker{\mathop{\mathrm{Coker}}}
\def\Coim{\mathop{\mathrm{Coim}}}

% Boxtimes
%
\DeclareMathSymbol{\boxtimes}{\mathbin}{AMSa}{"02}

%
% Macros for moduli stacks/spaces
%
\def\QCohstack{\mathcal{QC}\!\mathit{oh}}
\def\Cohstack{\mathcal{C}\!\mathit{oh}}
\def\Spacesstack{\mathcal{S}\!\mathit{paces}}
\def\Quotfunctor{\mathrm{Quot}}
\def\Hilbfunctor{\mathrm{Hilb}}
\def\Curvesstack{\mathcal{C}\!\mathit{urves}}
\def\Polarizedstack{\mathcal{P}\!\mathit{olarized}}
\def\Complexesstack{\mathcal{C}\!\mathit{omplexes}}
% \Pic is the operator that assigns to X its picard group, usage \Pic(X)
% \Picardstack_{X/B} denotes the Picard stack of X over B
% \Picardfunctor_{X/B} denotes the Picard functor of X over B
\def\Pic{\mathop{\mathrm{Pic}}\nolimits}
\def\Picardstack{\mathcal{P}\!\mathit{ic}}
\def\Picardfunctor{\mathrm{Pic}}
\def\Deformationcategory{\mathcal{D}\!\mathit{ef}}
\usepackage{cleveref}
\crefname{originallemma}{original lemma}{original lemmas}
\Crefname{originallemma}{Original lemma}{Original lemmas}
\makeatletter
\newlabel{lemma-colimit-abelian-category}{{010D}{source}{}{sourcepointer-1}{}}
\newlabel{lemma-colimit-abelian-category@cref}{{[originallemma][0][]010D}{[1][0][]source}{}{}{}}
\makeatother
\crefname{originallemma}{original lemma}{original lemmas}
\Crefname{originallemma}{Original lemma}{Original lemmas}
\makeatletter
\newlabel{categories-lemma-exact-adjoint}{{0039}{source}{}{sourcepointer-2}{}}
\newlabel{categories-lemma-exact-adjoint@cref}{{[originallemma][0][]0039}{[1][0][]source}{}{}{}}
\makeatother
\title{A reflective subcategory of an abelian category is abelian}
\author{The Stacks Project Authors}
\date{stacks-project:homology.tex:local-snapshot-20261012; GNU FDL 1.2 or later}
\begin{document}
\maketitle
\noindent\textit{Editorial note.} All delivered mathematical text is the author's own, copied from the pinned Stacks Project chapter file homology.tex; no statement or proof step was rewritten, reordered or removed.\par\medskip
\noindent\textit{Editorial note.} The proof is the author's complete argument; every step the author writes out is inside the counted span and no closing sentence is a placeholder.\par\medskip
\noindent\textit{Editorial note.} Every labelled reference occurring inside the counted spans is carried in assembly.external\_references as a non-load-bearing navigation pointer (same-chapter ones at their definition site, cross-chapter ones at the invocation inside the counted span), so no load-bearing mathematics is replaced by a stub. No bibliography citation occurs inside the counted spans.\par\medskip
\noindent\textit{Editorial note.} The same-name check for this card scanned every staged Stacks chapter of this repository -- all sources/stacks-*/pinned-original-source.tex files, the card's own chapter excluded -- and the labels found are recorded in dedup.same\_name\_cross\_chapter on the card.\par\medskip
\noindent\textit{Editorial note.} The preamble is the chapter's own native preamble excerpt already used for this source: the comment-type reference, slogan and history environments (native preamble lines 8-13), the Xy-pic setup (native preamble lines 14-18) and the native theorem declarations with the mathematical macros (native preamble lines 154-217, the actual last line of the file).\par\medskip
\noindent\textit{Allowed standard prerequisites (editorial).}\begin{itemize}
\item Existence of finite limits and colimits in an abelian category. Same chapter, Lemma lemma-colimit-abelian-category. Supplies the existence of finite limits and colimits and the construction of kernels and cokernels from equalizers and coequalizers.
\item Exactness of adjoint functors. Categories, Lemma categories-lemma-exact-adjoint. Shows that the left exact left adjoint is exact and that the right adjoint is left exact.
\end{itemize}\medskip
\ifdefined\hypertarget\hypertarget{sourcepointer-1}{}\fi
\noindent\textit{Source reference.} Original lemma 010D: sources/stacks-homology/pinned-original-source.tex, source lines 902--902. Navigation pointer invoked at proof line 1893; the counted span contains the pointer only and the cited result is not reproduced inside it, so the pointer is registered as a non-load-bearing navigation stub in the established form for Stacks material. No bibliography citation is involved.\par\medskip
\ifdefined\hypertarget\hypertarget{sourcepointer-2}{}\fi
\noindent\textit{Source reference.} Original lemma 0039: sources/stacks-categories/pinned-original-source.tex, source lines 3727--3727. Navigation pointer invoked at proof line 1896; the counted span contains the pointer only and the cited result is not reproduced inside it, so the pointer is registered as a non-load-bearing navigation stub in the established form for Stacks material. No bibliography citation is involved.\par\medskip
\begin{lemma}
\label{lemma-adjoint-get-abelian}
Let $a : \mathcal{A} \to \mathcal{B}$ and $b : \mathcal{B} \to \mathcal{A}$
be functors. Assume that
\begin{enumerate}
\item $\mathcal{A}$, $\mathcal{B}$ are additive categories,
$a$, $b$ are additive functors, and $a$ is right adjoint to $b$,
\item $\mathcal{B}$ is abelian and $b$ is left exact, and
\item $ba \cong \text{id}_\mathcal{A}$.
\end{enumerate}
Then $\mathcal{A}$ is abelian.
\end{lemma}

\begin{proof}
As $\mathcal{B}$ is abelian we see that all finite limits and colimits
exist in $\mathcal{B}$ by Lemma \ref{lemma-colimit-abelian-category}.
Since $b$ is a left adjoint we see that $b$ is also right exact
and hence exact, see
Categories, Lemma \ref{categories-lemma-exact-adjoint}.
Let $\varphi : B_1 \to B_2$ be a morphism of $\mathcal{B}$.
In particular, if $K = \Ker(B_1 \to B_2)$, then $K$ is
the equalizer of $0$ and $\varphi$ and hence
$bK$ is the equalizer of $0$ and $b\varphi$, hence
$bK$ is the kernel of $b\varphi$. Similarly, if
$Q = \Coker(B_1 \to B_2)$, then $Q$ is
the coequalizer of $0$ and $\varphi$ and hence
$bQ$ is the coequalizer of $0$ and $b\varphi$, hence
$bQ$ is the cokernel of $b\varphi$. Thus we see that every morphism
of the form $b\varphi$ in $\mathcal{A}$ has a kernel and a cokernel.
However, since $ba \cong \text{id}$ we see that every morphism of
$\mathcal{A}$ is of this form, and we conclude that kernels and
cokernels exist in $\mathcal{A}$. In fact, the argument shows that
if $\psi : A_1 \to A_2$ is a morphism then
$$
\Ker(\psi) = b\Ker(a\psi),
\quad\text{and}\quad
\Coker(\psi) = b\Coker(a\psi).
$$
Now we still have to show that $\Coim(\psi)= \Im(\psi)$.
We do this as follows.
First note that since $\mathcal{A}$ has kernels and cokernels it
has all finite limits and colimits (see proof of
Lemma \ref{lemma-colimit-abelian-category}).
Hence we see by Categories, Lemma \ref{categories-lemma-exact-adjoint}
that $a$ is left exact and
hence transforms kernels (=equalizers) into kernels.
\begin{align*}
\Coim(\psi)
& =
\Coker(\Ker(\psi) \to A_1)
& \text{by definition} \\
& =
b\Coker(a(\Ker(\psi) \to A_1))
& \text{by formula above} \\
& =
b\Coker(\Ker(a\psi) \to aA_1))
& a\text{ preserves kernels} \\
& =
b\Coim(a\psi)
& \text{by definition} \\
& =
b\Im(a\psi)
& \mathcal{B}\text{ is abelian} \\
& =
b\Ker(aA_2 \to \Coker(a\psi))
& \text{by definition} \\
& =
\Ker(baA_2 \to b\Coker(a\psi))
& b\text{ preserves kernels} \\
& =
\Ker(A_2 \to b\Coker(a\psi))
& ba = \text{id}_\mathcal{A} \\
& =
\Ker(A_2 \to \Coker(\psi))
& \text{by formula above} \\
& =
\Im(\psi)
& \text{by definition}
\end{align*}
Thus the lemma holds.
\end{proof}

\end{document}
